\section{Numerical and formal verification scope}\label{si:verification}

\subsection{Analytic results and finite checks}

The analytic claims in the main text and SI~S2--S4 are supported by the proofs supplied there. Finite calculations check specified cases; they do not establish continuum statements.

The exact-integer verifier cited in the main text enumerates $29{,}791$ states $(x,s_1,s_2)$ on a scale-40 grid and records 25 checks of witness classifications, region identities, licensing thresholds, endpoint comparisons, and hold-prefix persistence (Abaee, 2026d). This is exact arithmetic on a finite grid, not a proof over the continuum. The cited benchmark software reports 24 rational checks for the declared abstract action tubes and associated model quantities (Abaee, 2026f); it does not supply empirical income dynamics or calibrate a fishery model.

The diagonal CES thresholds and the pointwise off-diagonal geometric and harmonic criteria are proved in SI~S2--S3. They do not establish a global fixed-sum stationary-point structure, equality of off-diagonal minima, monotonicity of the critical elasticity, or a global phase diagram. No such global claim is made.

\subsection{Formalization scope}

A Lean source file for selected algebraic witness identities and CES rung examples is available at the cited immutable source revision: \href{https://github.com/MIKEAA2020/general-sustainability/blob/6f626ee48b2a3a59282b5df5adda76c8d5993fdc/lean/Formalizations/P1_AssessmentSeparation.lean}{formalization source}. It does not formalize the general real-exponent nesting, continuity/IVT transfer, or relative-interior argument. Those steps are supported by the written mathematical proofs and are not claimed as machine-checked results.

\subsection{Data and code availability}

No new empirical dataset was generated. The Northern-cod values are taken from DFO (2016), Table~A2: the reported 1983--1989 mean is $884.58$\,kt, rounded to $884.6$\,kt, and the 2015 value is approximately $298.65$\,kt. They are used only as a historical dimensional comparison, not as a current threshold. The cited code records provide the finite checks described above but do not constitute a complete reproduction of every calculation reported here; all analytic claims are supported by the proofs in the article and Supplementary Information.

\section{Literature positioning}\label{si:positioning}

\paragraph{Commensurability} The foundation statement that values may be only weakly comparable -- and that this incommensurability is constitutive of ecological economics -- is due to Martinez-Alier, Munda and O'Neill (1998). The main article's operators give one formal model of the distinction that paper draws between strong and weak commensurability.

\paragraph{Maximin and indicator formulations} The characterization of sustainability indicators and the role of maximin and MSY-type reasoning in operationalizing strong-sustainability criteria are developed in Martinet (2011), Cairns and Martinet (2014), and Doyen and Gajardo (2020). This literature includes formal viability and Pareto formulations; the typed operator is one narrow transition-certification idealization and is not presented as a comprehensive or uniquely strongest formalization.

\paragraph{Compensability in multi-criteria decision analysis} Compensability analysis is established in MCDA, including mappings between compensatory aggregation, weak sustainability, and outranking methods (Cinelli, Coles and Kirwan, 2014; Sch\"ar, Pohl and Geldermann, 2025). Weighted-sum limitations on nonconvex Pareto fronts are also well known (Das and Dennis, 1997). The present finite-menu theorem shares convex-hull geometry with these results; it asks a different question about state acceptance and plan selection, and the same quantifier gap can occur for static alternatives.

\paragraph{Acceptance semantics and decision aiding} The protocols of main-text Table~\ref{main-tab:protocols} are decision-aiding constructs: each specifies the rule by which a composite reading sorts a transition into the accepted or rejected category -- the sorting problematic of decision aiding (Roy, 1996; Vincke, 1992). Under the weight-adaptive rule a transition is accepted when every weighting licenses some plan, possibly a different plan for each weighting; the licensing thresholds are then the verdict's weight-sensitivity intervals, analogous to sensitivity analyses applied to composite-indicator rankings (Nardo et al., 2008). The preference-elicitation background is multi-attribute value theory's substitution rates (Keeney and Raiffa, 1976) and the incommensurable-criteria treatment of sustainable development (Munda, 2005).

\paragraph{Viability and robust control} The backward recursion is a typed instance of robust-predecessor, reach--avoid, capture-basin and hybrid-reachability constructions (Aubin, 1991; Aubin, Bayen and Saint-Pierre, 2011; Saint-Pierre, 1994; Lygeros, Tomlin and Sastry, 1999; Aubin and Catt\'e, 2002), with natural-resource management as an established application (De Lara and Doyen, 2008; Doyen et al., 2012). Fractional physical blending and complete-excursion concatenation compare distinct control classes; they are not a mixed/behavioral strategy equivalence. The here-and-now/wait-and-see analogy from adjustable robust optimization concerns the order of plan and weight quantifiers only; the weight is not modeled as a realized disturbance (Ben-Tal et al., 2004).

\paragraph{Planetary boundaries and the safe-and-just space} The typed operator is a formal idealisation of the separate-boundary reading underlying the Planetary Boundaries framework (Rockstr\"om et al., 2009; Steffen et al., 2015), and the same separately-checked-floor logic underlies the safe-and-just-space literature (Raworth, 2012; O'Neill et al., 2018; Fanning et al., 2022). Boundary processes interact and amplify human impacts (Lade et al., 2020), which is why the coupled disturbance regime of S4 is not a remote possibility. The main theorem provides a formal model of \emph{one reason} a noncompensatory stance may be adopted; it does not prove that planetary boundaries ought to be noncompensatory, which is a normative and scientific-design question.

\section{Finite-generation consequence under compact confinement}\label{si:finitegeneration}

A compactness consequence follows for nested generation constraints. Let $K$ be a compact state set, and let $V^{(k)}\subseteq K$ be closed, nested decreasing generation constraints with $\bigcap_{k\ge1}V^{(k)}=\varnothing$.

\begin{proposition}[finite-generation obstruction]\label{prop:finitegeneration}
Under these assumptions there is a finite generation $N$ for which $V^{(N)}=\varnothing$. Consequently, no trajectory confined to $K$ can satisfy the requirements through every generation.
\end{proposition}
\begin{proof}
If every $V^{(k)}$ were nonempty, the nested family of closed subsets of compact $K$ would have the finite-intersection property. Compactness would imply $\bigcap_{k\ge1}V^{(k)}\ne\varnothing$, contrary to the hypothesis. Hence some $V^{(N)}$ is empty. Since the sets are decreasing, every later generation set is also empty.
\end{proof}

The common compact confinement is essential to this finite-generation conclusion. Without it, take constant dynamics on $\mathbb R$ and $V^{(k)}=[k,\infty)$. Every finite prefix through generation $N$ is feasible from $x=N$, but no state satisfies all generations. The initial state used for each finite prefix escapes every common compact set.

For fixed dynamics and a fixed policy class, if safe sets satisfy $K_1\subseteq K_2$, then the viability kernel obeys $\operatorname{Viab}(K_1)\subseteq\operatorname{Viab}(K_2)$. Thus enlarging the safe set cannot contract the viability kernel under these assumptions. A different result would require a change in the dynamics, objective, or policy class.

\section*{References cited in the Supplementary Information}
\begin{itemize}[leftmargin=1.5em,itemsep=1pt]
\item Abaee, A. (2026a). An obstruction calculus for viability under incomplete observation. Zenodo. \url{https://doi.org/10.5281/zenodo.22552616}.
\item Abaee, A. (2026c). Does a surplus-production ladder improve forecasts of Northern cod? A scored test on NAFO 2J3KL. Zenodo. \url{https://doi.org/10.5281/zenodo.22553609}.
\item Abaee, A. (2026d). Verification code and figure pipeline for ``Aggregate indices and transition safety''. Figshare, version 3. \url{https://doi.org/10.6084/m9.figshare.33764023.v3}.
\item Abaee, A. (2026e). Delay-induced regime change in harvested stocks: the mobilising and protective channels of institutional feedback, and the review interval as control. Zenodo. \url{https://doi.org/10.5281/zenodo.22554217}.
\item Abaee, A. (2026f). \emph{SafeTransition: exact rational certification of transition safety for sustainability assessment}. Software release 1.3.0, Figshare, version 1. \url{https://doi.org/10.6084/m9.figshare.33971389.v1}.
\item Aubin, J.-P. (1991). \emph{Viability Theory}. Birkh\"auser, Boston. Aubin, J.-P., Bayen, A.~M., and Saint-Pierre, P. (2011). \emph{Viability Theory: New Directions}, 2nd ed. Birkh\"auser. Aubin, J.-P., and Catt\'e, F. (2002). \emph{Set-Valued Analysis}, 10(4), 379--416.
\item Ben-Tal, A., Goryashko, A., Guslitzer, E., and Nemirovski, A. (2004). \emph{Mathematical Programming}, 99(2), 351--376.
\item Cairns, R.~D., and Martinet, V. (2014). \emph{European Economic Review}, 69, 4--17.
\item Cardaliaguet, P. (1996). \emph{SIAM Journal on Control and Optimization}, 34(4), 1441--1460.
\item Cinelli, M., Coles, S.~R., and Kirwan, K. (2014). \emph{Ecological Indicators}, 46, 138--148.
\item Das, I., and Dennis, J.~E. (1997). \emph{Structural Optimization}, 14, 63--69.
\item Dearing, J.~A., et al. (2014). Safe and just operating spaces for regional social-ecological systems. \emph{Global Environmental Change}, 28, 227--238.
\item De Lara, M., and Doyen, L. (2008). \emph{Sustainable Management of Natural Resources}. Springer. Doyen, L., et al. (2012). \emph{Ecological Economics}, 75, 32--42. Doyen, L., and Gajardo, P. (2020). \emph{Natural Resource Modeling}, 33(3), e12250.
\item DFO (2016). \emph{Stock Assessment of Northern Cod (NAFO Divs.\ 2J3KL) in 2016}. CSAS Science Advisory Report 2016/026.
\item Fanning, A.~L., O'Neill, D.~W., Hickel, J., and Roux, N. (2022). \emph{Nature Sustainability}, 5(1), 26--36.
\item Keeney, R.~L., and Raiffa, H. (1976). \emph{Decisions with Multiple Objectives}. Wiley.
\item Lade, S.~J., et al. (2020). \emph{Nature Sustainability}, 3(2), 119--128.
\item Lygeros, J., Tomlin, C., and Sastry, S. (1999). \emph{Automatica}, 35(3), 349--370.
\item Martinet, V. (2011). \emph{Journal of Environmental Economics and Management}, 61(2), 183--197.
\item Martinez-Alier, J., Munda, G., and O'Neill, J. (1998). \emph{Ecological Economics}, 26(3), 277--286.
\item Munda, G. (2005). In \emph{Multiple Criteria Decision Analysis: State of the Art Surveys}. Springer, 953--986.
\item Nardo, M., et al. (2008). \emph{Handbook on Constructing Composite Indicators}. OECD.
\item O'Neill, D.~W., Fanning, A.~L., Lamb, W.~F., and Steinberger, J.~K. (2018). \emph{Nature Sustainability}, 1(2), 88--95.
\item Raworth, K. (2012). \emph{Oxfam Policy and Practice: Climate Change and Resilience}, 8(1), 1--26.
\item Rockstr\"om, J., et al. (2009). \emph{Nature}, 461, 472--475. Steffen, W., et al. (2015). Planetary boundaries: Guiding human development on a changing planet. \emph{Science}, 347(6223), 1259855.
\item Roy, B. (1996). \emph{Multicriteria Methodology for Decision Aiding}. Kluwer. Vincke, P. (1992). \emph{Multicriteria Decision-Aid}. Wiley.
\item Saint-Pierre, P. (1994). \emph{Applied Mathematics and Optimization}, 29, 187--209.
\item Schaefer, M.~B. (1954). \emph{Inter-American Tropical Tuna Commission Bulletin}, 1(2), 23--56.
\item Sch\"ar, S., Pohl, E., and Geldermann, J. (2025). \emph{Journal of Multi-Criteria Decision Analysis}, 32, e70013.
\item Sion, M. (1958). \emph{Pacific Journal of Mathematics}, 8(1), 171--176. von Neumann, J. (1928). \emph{Mathematische Annalen}, 100, 295--320.
\end{itemize}

\end{document}
